\section{Incidente 1: la retroalimentación positiva entre Retry, Cron y Rebuild}

El primer incidente de la clase no fue «se cayó una base de datos», sino la respuesta simultánea de varios mecanismos automáticos: el timeout dispara retry, el repair job y el cron aumentan la carga, la queue se alarga, y los cache miss junto con el rebuild terminan de saturar las dependencias. La primera acción de Incident response no es agregar más trabajos de reparación, sino identificar la retroalimentación positiva y detener al producer.

\subsection{Del initial fault a la cascade}

Esta sección reescribe el incidente de la clase como un dependency feedback graph, en lugar de relatar las operaciones en orden cronológico. Al leer la figura hay que seguir el recorrido desde el initial fault hasta la automatic reaction, la load amplification y la mitigation: lo decisivo no es qué componente reportó el error primero, sino qué retry, repair job y rebuild siguen aumentando la entrada de la saturated dependency, y qué punto de control puede reducir más rápido la ganancia del sistema.

\lecturefigure{07-incident-cascade.png}{Los retry y los repair job automáticos pueden amplificar una falla local hasta convertirla en un incidente global.}{Subtítulos locales 19:10--25:30; redibujo conceptual.}

\begin{knowledgebox}{Lectura de la figura: el orden de recuperación es el inverso del orden de optimización habitual}
En condiciones normales se busca throughput y automatización; durante un incidente primero se protege el invariant: detener los producer no críticos, limitar los retry, aplicar shed load y congelar los deploy, y después restablecer la read/write path mínima. Solo cuando la dependency tiene headroom se liberan gradualmente el backlog y el rebuild.
\end{knowledgebox}

\begin{importantbox}{Retry amplification}
Si la tasa original de solicitudes es $\lambda$, cada falla dispara en promedio $r$ reintentos y la probabilidad de falla es $p$, la carga adicional puede aproximarse como
\[
\lambda_{effective}\approx \lambda(1+pr+(pr)^2+\cdots).
\]
Cuando $pr$ se acerca a 1, la cadena de reintentos se vuelve muy sensible. Los sistemas reales controlan la amplificación con bounded retry, exponential backoff, jitter, circuit breaker y un global retry budget.
\end{importantbox}

\teachervoice{El relato del incidente en la clase tiene un fuerte tono dramático de «una persona sentada ahí salvando el sistema». Lo que de verdad hay que aprender es la command structure: quién puede detener al producer, quién evalúa el invariant de los datos, quién se encarga de la comunicación con los usuarios y qué automatizaciones deben desactivarse en degraded mode.}

\subsection{Resumen de la sección}

La causa raíz de una Incident cascade suele ser un feedback loop, no un componente aislado. La recuperación exige primero reducir la ganancia del sistema y después reparar el estado; el postmortem debe incorporar retry, cron, queue y migration en un mismo dependency graph.
